\documentclass[a4paper]{article}

\usepackage[spanish]{babel}
\usepackage{graphicx}
\usepackage{amsmath, amssymb, amsthm}
\usepackage[margin=2cm]{geometry}
\usepackage{fancyhdr}
\usepackage{enumerate}
\usepackage[shortlabels]{enumitem}
\usepackage{parskip}
\usepackage[hidelinks]{hyperref}
\usepackage{float}
\usepackage{booktabs}
\usepackage{array}



\renewcommand{\proofname}{Demostración}

\newtheorem{teorema}{Teorema}
\newtheorem{lema}{Lema}
\newtheorem{definicion}{Definición}
\newtheorem{corolario}{Corolario}



% cabecera
\pagestyle{fancy}
\fancyhead[l]{Mecánica Celeste}
\fancyhead[c]{Grafos}
\fancyhead[r]{\today}
\fancyfoot[c]{\thepage}
\renewcommand{\headrulewidth}{0.2pt} % linea horizontal

\begin{document}

\begin{titlepage}
  \centering
  {\Large Universidad de Antioquia\par}
  \vspace{2cm}
  {\huge\bfseries Mecánica Celeste\par}
  \vspace{0.4cm}
  {\Large Grafos en sistemas jerárquicos \par}
  \vspace{2.5cm}
  {\large Autor:\par}
  \vspace{0.2cm}
  {\large Daniel Soto Villada\par}
  {\large Estudiante de Astronomía\par}
  \vfill
  {\large \today\par}
\end{titlepage}

\newpage
\tableofcontents
\newpage


Los sistemas estelares múltiples jerárquicos son, en el sentido
matemático preciso, árboles binarios plenos con raíz y pesados.
La teoría de grafos no solo describe la estructura: también formaliza
exactamente por qué la jerarquía garantiza estabilidad dinámica.

% -----------------------------------------------------------------------
\section{Fundamentos matemáticos}

\begin{definicion}[Grafo]
Un \textbf{grafo} es un par $G = (V, E)$ donde $V$ es un conjunto finito
de vértices y
$$E \subseteq \bigl\{\{u,v\} : u,v \in V,\; u \neq v\bigr\}$$
es el conjunto de aristas.
\end{definicion}

Dos propiedades básicas se derivan inmediatamente:

\begin{itemize}
  \item \textbf{Grado de un vértice} $v$:
        $\deg(v) = \bigl|\{u \in V : \{u,v\} \in E\}\bigr|$
  \item \textbf{Lema del apretón de manos}:
        $\displaystyle\sum_{v \in V} \deg(v) = 2|E|$
\end{itemize}

\begin{definicion}[Árbol]
Un árbol $T = (V, E)$ es un grafo conexo y acíclico. Son equivalentes:
\begin{enumerate}
  \item $T$ es un árbol.
  \item $|E| = |V| - 1$ y $T$ es conexo.
  \item Existe exactamente un camino entre cualquier par $u, v \in V$.
\end{enumerate}
\end{definicion}

\begin{definicion}[Árbol con raíz]
Un árbol con raíz $(T, r)$ distingue un vértice $r \in V$ llamado
raíz, lo que induce:
\begin{itemize}
  \item la \textbf{profundidad} $\operatorname{prof}(v)$, definida como
        la longitud del camino de $r$ a $v$,
  \item la función $\operatorname{padre}(v)$ para todo $v \neq r$,
  \item el conjunto $\operatorname{hijos}(v)$,
  \item la distinción entre \textbf{hojas}
        $L = \{v : \operatorname{hijos}(v) = \varnothing\}$
        y \textbf{nodos internos} $I = V \setminus L$.
\end{itemize}
\end{definicion}

\begin{definicion}[Árbol binario pleno]
Un árbol binario pleno (\textit{full binary tree}) es aquel en que
todo nodo interno tiene exactamente $2$ hijos.
\end{definicion}

\begin{teorema}
Sea $T$ un árbol binario pleno con $k$ hojas. Entonces:
$$|I| = k - 1, \qquad |V| = 2k - 1, \qquad |E| = 2k - 2.$$
\end{teorema}

\begin{proof}
Por inducción sobre $k$. El caso base $k = 1$ da $|I| = 0$, $|V| = 1$,
$|E| = 0$, que satisface las tres igualdades. Para el paso inductivo,
si una hoja se convierte en nodo interno al ganar $2$ hijos, el número
neto de hojas aumenta en $1$ y el de nodos internos también en $1$,
preservando la relación $|I| = k - 1$.
\end{proof}

% -----------------------------------------------------------------------
\section{Castor como árbol con raíz}

El sistema se formaliza como $(T, \Omega)$, donde $\Omega$ es el
baricentro global:

$$V = \{\Omega,\;\mathrm{CM_{AB}},\;\mathrm{CM_A},\;\mathrm{CM_B},\;
       \mathrm{CM_C},\;\mathrm{Aa},\;\mathrm{Ab},\;\mathrm{Ba},\;
       \mathrm{Bb},\;\mathrm{Ca},\;\mathrm{Cb}\}$$

\begin{align*}
E = \bigl\{
  &\{\Omega, \mathrm{CM_{AB}}\},\;
   \{\Omega, \mathrm{CM_C}\},\;
   \{\mathrm{CM_{AB}}, \mathrm{CM_A}\},\;
   \{\mathrm{CM_{AB}}, \mathrm{CM_B}\},\\
  &\{\mathrm{CM_A}, \mathrm{Aa}\},\;
   \{\mathrm{CM_A}, \mathrm{Ab}\},\;
   \{\mathrm{CM_B}, \mathrm{Ba}\},\;
   \{\mathrm{CM_B}, \mathrm{Bb}\},\\
  &\{\mathrm{CM_C}, \mathrm{Ca}\},\;
   \{\mathrm{CM_C}, \mathrm{Cb}\}
\bigr\}
\end{align*}

Propiedades inmediatas:
$$|V| = 11, \qquad |E| = 10 = 11 - 1 \;\checkmark$$

Con $6$ hojas (estrellas reales) y $5$ nodos internos (baricentros),
se verifica $|I| = 6 - 1 = 5\;\checkmark$, confirmando que el árbol
es \textbf{binario pleno}.

Verificación del lema del apretón de manos:
\begin{align*}
\sum_{v \in V} \deg(v)
  &= \deg(\Omega) + \deg(\mathrm{CM_{AB}}) + \deg(\mathrm{CM_A})
   + \deg(\mathrm{CM_B}) + \deg(\mathrm{CM_C})
   + \sum_{\text{hojas}} \deg(v)\\
  &= 2 + 3 + 3 + 3 + 3 + (6 \times 1) = 20 = 2 \times 10 \;\checkmark
\end{align*}

Las profundidades revelan que el árbol es binario pleno pero
\textbf{no completo ni perfecto}:
$$\operatorname{prof}(\mathrm{Ca}) = \operatorname{prof}(\mathrm{Cb}) = 2,
\qquad
\operatorname{prof}(\mathrm{Aa}) = \operatorname{prof}(\mathrm{Ab})
= \operatorname{prof}(\mathrm{Ba}) = \operatorname{prof}(\mathrm{Bb}) = 3.$$

% -----------------------------------------------------------------------
\section{Árbol pesado y la condición de estabilidad}

\begin{definicion}[Grafo pesado]
Un grafo pesado $G = (V, E, w)$ agrega una función de peso
$w : E \to \mathbb{R}^+$. Para sistemas estelares, el peso natural es
el semieje mayor orbital de cada órbita:
$$w\bigl(\{u,v\}\bigr) = a_{u,v}.$$
\end{definicion}

La estabilidad jerárquica se traduce en una condición sobre los
cocientes de pesos. Para cada nodo interno $v \neq \Omega$ se define el
\textbf{cociente de separación}:
$$\rho(v) = \frac{w\bigl(\{\operatorname{padre}(v),\, v\}\bigr)}
                 {w\bigl(\{v,\, c\}\bigr)},
\qquad \forall\; c \in \operatorname{hijos}(v).$$

Un sistema es jerárquicamente estable si y solo si
$\rho(v) \gg 1$ para todos los nodos internos. En la práctica se
requiere $\rho \gtrsim 3\text{--}5$; para longevidad real, $\rho > 10$.

Para Castor (valores aproximados en UA):

\begin{center}
\begin{tabular}{lcc}
\toprule
Arista & Peso $a$ (UA) & $\rho$ \\
\midrule
$\mathrm{Aa}\text{--}\mathrm{Ab}$ & $\approx 0.09$ & --- \\
$\mathrm{Ba}\text{--}\mathrm{Bb}$ & $\approx 0.04$ & --- \\
$\mathrm{CM_A}\text{--}\mathrm{CM_B}$ & $\approx 73$
  & $\rho(\mathrm{CM_{AB}}) \approx 73/0.09 \approx 810$ \\
$\Omega\text{--}\mathrm{CM_C}$ & $\approx 1000$
  & $\rho(\Omega) \approx 1000/73 \approx 14$ \\
\bottomrule
\end{tabular}
\end{center}

Ambos cocientes satisfacen $\rho \gg 1$, lo que justifica la
estabilidad del sistema.

% -----------------------------------------------------------------------
\section{La notación de Newick como serialización del árbol}

La notación compacta \texttt{((Aa,Ab),(Ba,Bb))} es exactamente el
\textbf{formato Newick}, usado en filogenética desde 1974 para
serializar árboles filogenéticos. Formalmente es una función
recursiva:

$$\mathrm{newick}(v) =
\begin{cases}
  \operatorname{nombre}(v)
    & \text{si } v \text{ es hoja,}\\[4pt]
  \bigl(\mathrm{newick}(c_1),\;\mathrm{newick}(c_2)\bigr)
    & \text{si } v \text{ es nodo interno.}
\end{cases}$$

Para Castor:
$$\texttt{(((Aa,Ab),(Ba,Bb)),(Ca,Cb))}$$

Es una codificación biyectiva del árbol que mapea toda su
estructura a una cadena. Es notable que astrofísicos y filogenetistas
hayan convergido en la misma representación para describir
jerarquías anidadas en dominios completamente distintos.

% -----------------------------------------------------------------------
\section{Sobre el ejemplo de Castor}

La descripción presentada es correcta en todos sus aspectos: la
clasificación séxtuple, la jerarquía de tres niveles, el árbol
y la notación compacta corresponden fielmente al sistema
$\alpha$~Geminorum. Si hay algo que matizar: la separación entre el
par~AB y YY~Gem es de $\sim 1000\;\mathrm{UA}$ con un período
orbital estimado en varios miles a decenas de miles de años, lo que
hace difícil confirmar la órbita directamente; sin embargo, el
movimiento propio compartido confirma la ligadura gravitacional. No es
un error en la descripción, sino una limitación observacional del
nivel más externo del árbol.

% -----------------------------------------------------------------------
\section{Ejemplo inventado: Sistema Múltiple Ficticio Zeta Aquilae-X
         (séptuple)}

Se define un sistema con una asimetría más pronunciada que Castor:
la rama derecha contiene solo dos estrellas ligadas directamente al
baricentro intermedio, mientras que la rama izquierda anida tres
niveles de binarias. Además, $C$ es una \textbf{estrella individual}
(no binaria), lo que crea un nodo interno con un hijo hoja y un hijo
interno.

Notación de Newick:
$$\texttt{(((Aa,Ab),(Ba,Bb)),C),(Da,Db)}$$

Formalización:
\begin{align*}
V = \bigl\{&\Omega,\;\mathrm{BM_4},\;\mathrm{BM_5},\;\mathrm{BM_3},\;
             \mathrm{BM_1},\;\mathrm{BM_2},\\
           &C,\;\mathrm{Da},\;\mathrm{Db},\;
            \mathrm{Aa},\;\mathrm{Ab},\;\mathrm{Ba},\;\mathrm{Bb}
    \bigr\}
\end{align*}

\begin{align*}
E = \bigl\{
  &\{\Omega, \mathrm{BM_4}\},\;
   \{\Omega, \mathrm{BM_5}\},\;
   \{\mathrm{BM_4}, \mathrm{BM_3}\},\;
   \{\mathrm{BM_4}, C\},\;
   \{\mathrm{BM_5}, \mathrm{Da}\},\\
  &\{\mathrm{BM_5}, \mathrm{Db}\},\;
   \{\mathrm{BM_3}, \mathrm{BM_1}\},\;
   \{\mathrm{BM_3}, \mathrm{BM_2}\},\;
   \{\mathrm{BM_1}, \mathrm{Aa}\},\\
  &\{\mathrm{BM_1}, \mathrm{Ab}\},\;
   \{\mathrm{BM_2}, \mathrm{Ba}\},\;
   \{\mathrm{BM_2}, \mathrm{Bb}\}
\bigr\}
\end{align*}

Propiedades:
\begin{align*}
|V| &= 13, \qquad |E| = 12 = 13 - 1 \;\checkmark\\
|L| &= |\{\mathrm{Aa},\mathrm{Ab},\mathrm{Ba},\mathrm{Bb},C,
           \mathrm{Da},\mathrm{Db}\}| = 7\\
|I| &= 6 = 7 - 1 \;\checkmark \quad \Rightarrow \text{árbol binario pleno}
\end{align*}

Las profundidades de las hojas son:
$$\operatorname{prof}(C)
= \operatorname{prof}(\mathrm{Da})
= \operatorname{prof}(\mathrm{Db}) = 2,
\qquad
\operatorname{prof}(\mathrm{Aa})
= \operatorname{prof}(\mathrm{Ab})
= \operatorname{prof}(\mathrm{Ba})
= \operatorname{prof}(\mathrm{Bb}) = 4.$$

La profundidad máxima es $4$, un nivel más que Castor. La secuencia
de grados:
\begin{align*}
\deg(\Omega) &= 2,\\
\deg(\mathrm{BM_4}) = \deg(\mathrm{BM_5}) = \deg(\mathrm{BM_3})
  = \deg(\mathrm{BM_1}) = \deg(\mathrm{BM_2}) &= 3,\\
\deg(C) = \deg(\mathrm{Da}) = \deg(\mathrm{Db})
  = \deg(\mathrm{Aa}) = \deg(\mathrm{Ab})
  = \deg(\mathrm{Ba}) = \deg(\mathrm{Bb}) &= 1.
\end{align*}

Verificación:
$$\sum_{v \in V} \deg(v) = 2 + (5 \times 3) + (7 \times 1)
= 2 + 15 + 7 = 24 = 2 \times 12 \;\checkmark$$

Lo que hace interesante este árbol desde la teoría de grafos es su
\textbf{asimetría estructural}: la rama izquierda de $\Omega$ alcanza
profundidad máxima $4$, mientras que la rama derecha solo llega a
profundidad $2$. La estrella $C$ es un caso especial: es hija directa
de $\mathrm{BM_4}$ sin formar una binaria propia, lo que significa que
$\mathrm{BM_4}$ tiene un hijo hoja y un hijo interno.

Finalmente, la condición de estabilidad exige que los pesos de las
aristas crezcan estrictamente de hoja a raíz a lo largo de cualquier
camino. Formalmente, para todo camino
$v_0, v_1, \ldots, v_\ell = \Omega$ con $v_0 \in L$:
$$w\bigl(\{v_0, v_1\}\bigr)
< w\bigl(\{v_1, v_2\}\bigr)
< \cdots
< w\bigl(\{v_{\ell-1}, \Omega\}\bigr).$$
Un árbol pesado con esa propiedad es un \textbf{árbol de Hasse
ponderado} y la condición de estabilidad es simplemente que ese
árbol sea estrictamente creciente en pesos de hoja a raíz.



\bibliographystyle{plainurl}
\bibliography{bibliografia} 
\end{document}
